\scanpage{3}\mspage{\unsure{1}}
1. Ens. à quatre éléments.

Un ens. à quatre éléments $I$ est muni canoniquement d'une structure \deleted{\unsure{de torseur sous un}} \inserted{d'espace affine de} dim $2$ sur $\mathbb{F}_2$, i.e. d'un sous-groupe $V\subset\unsure{\mathfrak{S}|G}_I$, isomorphe (non can!) à $\mathbb{F}_2^2$, tel que l'action de $V$ sur \deleted{$\unsure{\mathbb{F}_2}$} $I$ en fasse un torseur. Le sous-groupe $V\subset\mathfrak{S}_I$ est le seul sous-groupe distingué abélien $\ne\{0\}$, \deleted{\unsure{c'est aussi}} il est son propre \deleted{normalisateur} \inserted{centralisateur} i.e. $\mathfrak{S}_I/V$ opère fidèlement sur $V$, en fait on a
\[
\mathfrak{S}_I/V\xrightarrow{\sim}\mathrm{Aut}_{\mathrm{gr}}(V)=\mathrm{Aut}_{\mathbb{F}_2\text{-vect}}(V)\xrightarrow{\sim}\deleted{\mathrm{Aut}}\mathfrak{S}_J
\]
où $J=V^{\unsure{\times|*}}$, la suite exacte
\[
1\longrightarrow V\longrightarrow\mathfrak{S}_I\longrightarrow\mathfrak{S}_J\longrightarrow1
\]
étant \deleted{\unsure{isomorphe|identique}} \inserted{identique} à la suite
\[
1\unsure{\text{---}|\longrightarrow}\underset{\text{translations}}{V}\longrightarrow\mathrm{Aut}_{\mathrm{aff}}(I)\longrightarrow\mathrm{Aut}_{\mathbb{F}_2}(V)\longrightarrow1.
\]

Les \inserted{trois} éléments de $J=V^{\times}$ correspondent aux trois partitions de type $(2.2)$ de $I$
% slanted
\margin{ou aux trois involutions \unsure{sur|en} $I$ sans pt fixe}
(on peut l'interpréter \inserted{ainsi,} que les éléments de $J$ correspondent aux \uline{droites} \inserted{vectorielles} de $V^{*}$, qui définissent \deleted{\unsure{des|les}} donc des \deleted{\illegible{1}} relations d'équivalence de $I$, vers les espaces affines quotients de dim $1$) -- les classes \unsure{intervenant|inversement} dans l'ensemble de ces \inserted{3} relations d'équivalence correspondent aux $6$ \deleted{\unsure{applications|paires}} \inserted{parties} à deux éléments de $I$) ou aussi \deleted{\unsure{aux}} $6$ droites affines de l'espace affine $I$, l'application
\[
\begin{aligned}
\mathcal{P}_2(I)&\longrightarrow\mathcal{P}_{\unsure{2,2|1,1}}(\unsure{I|J})\\
A&\longmapsto\{A,I\setminus A\}
\end{aligned}\tag*{1}
\]
% archive sheet number: 2
